% =============================================================
% Section 2: Preliminaries and Related Work
% =============================================================

\section{Preliminaries and Related Work}
\label{sec:related}

\subsection{Time Series Foundation Models}
\label{sec:rel_tsfm}

The past two years have produced a generation of models pretrained on large,
heterogeneous time series corpora with the goal of zero-shot or few-shot
transfer~\cite{jin2023large,liang2024foundation}.
These models differ along two design axes: how input data is divided into
tokens and what prediction target guides training.

\textbf{Patch-based encoder-decoder models} divide the context window into
fixed-length patches and learn representations via masked modelling.
MOMENT~\cite{goswami2024moment} applies this to sub-series segments,
supporting forecasting, classification, and anomaly detection.
Moirai~\cite{woo2024moirai} extends it with a multi-patch-size architecture
that accommodates variable-frequency series.

\textbf{Decoder-only autoregressive models} treat forecasting as next-patch
prediction.
TimesFM~\cite{das2024timesfm} trains a 500M-parameter model on a corpus
spanning Google Trends, Wikipedia pageviews, and synthetic series.
Timer~\cite{liu2024timer} follows the same autoregressive design at smaller
scale.

\textbf{Language-model-inspired models} convert continuous series values into
discrete tokens---treating numerical forecasting as a text-prediction problem---and
fine-tune sequence-to-sequence architectures.
Chronos~\cite{ansari2024chronos} builds on T5~\cite{raffel2020t5}, casting
forecasting as a sequence-to-sequence language task.
Lag-Llama~\cite{rasul2024lagllama} adapts LLaMA by appending lagged copies of
the input series as additional context features.
Concurrent work has also explored prompting general-purpose LLMs directly for
zero-shot forecasting~\cite{gruver2023llmts}, with mixed results that further illustrate the gap between raw performance and
what the model has actually learned to represent.

\subsection{The Performance Debate: Promise vs.\ Reality}
\label{sec:rel_debate}

The rapid proliferation of TSFMs has been accompanied by an equally rapid
accumulation of contradictory evidence about their practical utility.
Understanding this debate is essential context for our geometric analysis.

\textbf{The sceptical challenge.}
\citet{zeng2023transformers} delivered the most influential critique:
a one-layer linear model (DLinear) matches or outperforms every major
Transformer-based forecaster of its era on five standard long-horizon benchmarks,
including ETT and Weather.
The implication is stark: if a model with no representational capacity matches a
deeply stacked Transformer, either the Transformer is learning something irrelevant
to the task, or the benchmarks are not measuring what matters.
This result reframed the research agenda from ``which Transformer is best'' to
``does Transformer complexity help at all''.

\textbf{The zero-shot argument.}
Proponents of TSFMs respond that the comparison is unfair: DLinear is trained
per-dataset, while TSFMs operate zero-shot.
The Chronos~\cite{ansari2024chronos} and TimesFM~\cite{das2024timesfm}
evaluations both show zero-shot performance competitive with, and sometimes
exceeding, supervised per-dataset models---a capability that linear models
cannot replicate by definition.
Moirai~\cite{woo2024moirai} further demonstrates that a single foundation model
can handle variable frequencies and multivariate series within a unified
architecture.

\textbf{The unresolved tension.}
Despite this back-and-forth, several uncomfortable facts remain.
First, the zero-shot superiority of TSFMs is not universal: on standard
long-horizon benchmarks with sufficient training data, supervised patch-based
models like PatchTST~\cite{nie2023patchtst} remain highly competitive at a
fraction of the parameter count.
Second, TSFM evaluations are not standardised: reported results vary
significantly across evaluation protocols, preprocessing choices, and
prediction horizons~\cite{woo2024moirai}, making apples-to-apples comparison
difficult.
Third, and most importantly for our work: \emph{the debate is conducted
entirely at the output level}.
Neither proponents nor sceptics have systematically examined what these models
actually store in their internal representations.

\textbf{The missing question.}
A model that achieves competitive zero-shot MAE on ETTh1 could do so by
(a)~building a rich, transferable internal representation of temporal
patterns, or (b)~memorising statistical properties of the pretraining data
that happen to generalise.
These two mechanisms have radically different implications for when TSFMs
generalise, when they fail, and how they should be used.
Performance metrics cannot distinguish them.
Geometric analysis of representations can.
This is the contribution of the present paper.

\subsection{Geometric Analysis of Neural Representations}
\label{sec:rel_geometry}

The internal geometry of learned representations has been studied systematically
in vision and language, providing the methodological toolkit we adapt to
time series.

\textbf{Intrinsic dimensionality.}
\citet{ansuini2019intrinsic} showed that deep image classifiers exhibit a
characteristic layer-wise intrinsic dimensionality (ID) profile: expansion in
early layers, compression to a low-dimensional manifold in terminal layers, with
the compression ratio correlating with generalisation.
\citet{birdal2021intrinsic} extended this to language models, showing that low
intrinsic dimensionality of the fine-tuning landscape explains why large
pretrained models adapt efficiently with few parameters.
The TwoNN estimator~\cite{facco2017estimating} provides a non-parametric method
for estimating this dimension from nearest-neighbour distance ratios---part of
the geometric toolkit we adapt to time series representations.

\textbf{Representational geometry in language models.}
\citet{reif2019visualizing} demonstrated that BERT organises contextual
embeddings into geometrically separable clusters corresponding to syntactic and
semantic categories---a structural finding invisible to accuracy benchmarks alone.
\citet{raghu2017svcca} used SVCCA to show that lower layers stabilise early
in training while upper layers continue evolving, showing how representations form layer by layer and motivating the analysis
of what learned representations store.

\textbf{The time series gap.}
No equivalent geometric characterisation exists for time series models.
We do not know whether TSFMs organise their embedding spaces around temporal
patterns, whether their internal structure is stable under small input changes,
or whether proximity in embedding space predicts forecast reliability.
These are the questions LGF is designed to answer.

\subsection{Explainability for Time Series}
\label{sec:rel_xai}

Existing XAI approaches for time series assign importance scores to input
timesteps---via LIME, SHAP, or integrated gradients---or measure model
sensitivity to input perturbations~\cite{ismail2020benchmarking,
theissler2022explainable}.
Both classes operate at the input level and explain individual predictions.
Neither provides any account of what a model has learned to encode in its
representation space.

Linear probing~\cite{alain2017probing} bridges this gap: it tests whether a
target property can be read out of frozen intermediate representations using
a linear classifier, revealing what information is stored without the
confounds of end-to-end training.
In NLP, probing revealed that BERT encodes syntactic structure in intermediate
layers and semantic structure in upper layers~\cite{reif2019visualizing}.
We apply this approach to provide the first direct account of what temporal
information is stored in pretrained TSFM representations.
